% Part: first-order-logic
% Chapter: syntax-and-semantics
% Section: introduction

\documentclass[../../../include/open-logic-section]{subfiles}

\begin{document}

\olfileid{fol}{syn}{int}

\olsection{سريزه}

د لومړۍ درجې منطق د تيورۍ او فرا تيورۍ د جوړولو دپاره بايد اول د هغې د عبارتونو نحو او معنٰی پوهنه تعريف کړو۔ د لومړۍ درجې منطق عبارتونه ترمونه او !!{formula}s دي۔ ترمونه له !!{variable}s، !!{constant}s او !!{function}s جوړېږي۔ بيا !!^{formula}s له !!{predicate}s او ترمونو سره يوځاے جوړېږي (همدا تر ټولو کوچنۍ، «اتومي» !!{formula}s جوړوي)، او له اتومي !!{formula}s څخه د منطقي نښلوونکو او کمیت ټاکونکو په مرسته لا پېچلي فارمولونه جوړولے شو۔ د جوړېدو د قواعدو د ليکلو ډېرې بېلابېلې طريقې شته؛ موږ يوازې يوه ممکنه طريقه ورکوو۔ نور نظامونه به نورې نښې غوره کوي، د لومړنيو نښلوونکو نور سټونه ټاکي، قوسونه به په بله طريقه کاروي (يا به يې بيخي نۀ کاروي، لکه په تش په نامه پولنډۍ نښه ليکنه کښې)۔ خو د ټولو طريقو ګډه خبره دا ده چې د جوړېدو قواعد د ترمونو او !!{formula}s سټ \emph{په استقرايي ډول} تعريفوي۔ که کار په سمه توګه وشي، هر عبارت د جوړېدو د قواعدو له مخې اساساً يوازې په يوه طريقه منځ ته راتلے شي۔ هغه استقرايي تعريف چې داسې عبارتونه جوړوي چې \emph{يکتا لوستونتيا} لري، دا معنا لري چې دغو عبارتونو ته په هماغه طريقه، يعنې په استقرايي تعريف، مانا ورکولے شو۔

عبارتونو ته مانا ورکول د معنٰی پوهنې کار دے۔ په معنٰی پوهنه کښې مرکزي مفهوم په !!a{structure} کښې صدق دے۔ !!^a{structure} د ژبې بنسټيزو اجزاوو ته مانا ورکوي: !!a{domain} د څيزونو يو غيرتشه سټ دے۔ کمیت ټاکونکي داسې تفسيرېږي چې پر همدې ساحه ګرځي؛ !!{constant}s ته د ساحې غړي، !!{function}s ته له !!{domain} څخه همدغې ساحې ته تابعې، او !!{predicate}s ته پر !!{domain} اړيکې ګومارل کېږي۔ !!{domain} او بنسټيزو لغتونو ته شوې ګومارنې په ګډه !!a{structure} جوړوي۔ !!^{variable}s په !!{formula}s کښې راتلے شي، او د معنٰی پوهنې د ورکولو دپاره بايد هغو ته هم د !!{domain} !!{element}s ګومارل شي؛ همدا د متغير ګومارنه ده۔ په پاې کښې د صدق اړيکه دغه څيزونه يوځاے کوي۔ !!^a{formula} په !!a{structure}~$\Struct{M}$ کښې د متغير ګومارنې~$s$ په نسبت صدق کولے شي، چې $\Sat{M}{!A}[s]$ ليکل کېږي۔ دا اړيکه هم د~$!A$ پر جوړښت په استقرا تعريفېږي؛ د منطقي نښلوونکو د رښتياوالي جدولونه کاروي، څو مثلاً د $!A \land !B$ صدق د $!A$ او ~$!B$ د صدق (يا نۀ صدق) له مخې تعريف کړي۔ بيا معلومېږي چې که !!{formula}~$!A$، !!a{sentence} وي، يعنې ازاد متغيرونه ونۀ لري، د متغير ګومارنه بې اغېزې ده؛ نو د !!{sentence}s په اړه يوازې دا ويلے شو چې په !!{structure}s کښې صدق کوي (يا نۀ کوي)۔

د جملو دپاره د صدق اړيکې $\Sat{M}{!A}$ پر بنسټ بيا د اعتبار، استلزام او د صدق وړتيا بنسټيز معنايي مفهومونه تعريفولے شو۔ جمله هله معتبره ده، $\Entails !A$، چې په هر جوړښت کښې صدق وکړي۔ د !!{sentence}s يو سټ ئې هله لازموي، $\Gamma \Entails !A$، چې هر هغه !!{structure} چې د~$\Gamma$ ټولو !!{sentence}s ته صدق ورکوي، د~$!A$ صدق هم ورکړي۔ او د جملو يو سټ هله د صدق وړ دے چې کوم !!{structure} د هغه ټولو !!{sentence}s ته په يوه وخت کښې صدق ورکړي۔ ځکه چې !!{formula}s په استقرايي ډول تعريفېږي، او صدق بيا د !!{formula}s پر جوړښت په استقرا تعريفېږي، استقرا د خپلې معنٰی پوهنې د ځانګړتياوو د ثابتولو او د تعريف شويو معنايي مفهومونو د تړلو دپاره کارولے شو۔

\end{document}
